
\documentclass[10pt,a4paper]{article}

% Basic packages
\usepackage[utf8]{inputenc}
\usepackage[T1]{fontenc}
\usepackage{graphicx}
\usepackage{amsmath}
\usepackage{booktabs}
\usepackage{titlesec}

% Paragraph formatting
\setlength{\parindent}{0pt}
\setlength{\parskip}{4pt}

% Section formatting
\setlength{\emergencystretch}{3em}

\begin{document}

% Title page
\begin{titlepage}

\begin{center}

\vspace*{2cm}

{\Large\textbf{WORLD POPULATION TRENDS AND CO$_2$ EMISSIONS}}

\vspace{1.5cm}

{\large\textbf{PROJECT REPORT}}

\vspace{1.5cm}

\textbf{Submitted by:}\\
Diane MUTONI

\vspace{0.5cm}

\textbf{Student ID:}\\
20252MBI033

\vspace{0.5cm}

\textbf{Submission Date:}\\
6 November 2026

\vspace{0.5cm}

\textbf{Course:}\\
R Programming

\vfill

\end{center}

\end{titlepage}


\section{Abstract}

This study examines world population trends and their relationship with CO2 emissions using population and CO2 emissions datasets. Exploratory data analysis was conducted to examine the structure, completeness, duplicates and outliers in the population dataset. Population projections for 2030 were calculated using the exponential growth model. Data visualization was used to identify 10 most populous countries and examine their  CO2 emission trends. Correlation analysis was performed to investigate the relationship between 2022 population size and 2019 CO2 emissions per capita. The correlation coefficient was 0.00817, indicating an almost nonexistent linear relationship.
Furthermore, continental comparisons showed that Asia had the highest sum of country-level per-capita emissions (242), while South America had the lowest (29.4). The outcomes showed that population size alone does not explain differences in CO2 emissions per capita.


\section{Introduction}

Population growth and CO2 emissions are important factors in understanding global environmental challenges. The increase in population may influence energy consumption, industrial activities and high demand for natural resources. However, countries with larger populations do not necessarily have higher CO2 emissions per person. This project analyzes population trends and CO2 emission patterns using data management, exploratory data analysis, visualization and statistical techniques in R.


\section{Methods and Dataset Description}

Two datasets were used: the World Population Dataset, containing 234 observations and 17 variables, and the CO2 Emissions Around the World dataset, containing 215 observations and 35 variables. The population dataset included country names, population figures for different years, area, density, growth rate and continent. The CO2 dataset contained annual CO2emissions per capita from 1990 to 2019.
R programming was used for importing both datasets,examining them, data cleaning, manipulation and visualization. Exploratory data analysis(EDA) involved examining variable names, data types, head/tail datasets, duplicates, missing values and outliers. The exponential population projection formula was applied to estimate population in 2030. The dplyr, tidyr and ggplot2 packages were used for data manipulation and visualization. Correlation and scatter plots were used to examine the relationship between population and CO2 emissions.


\section{Results and Findings}

The exploratory analysis showed that the population
dataset contained 234 countries and 17 variables,
with no missing values or duplicated observations.
However, several quantitative variables contained
outliers, particularly population and area.

Population projections for 2030 showed differences
in expected growth across countries. For example,
Afghanistan was projected to reach approximately
50.52 million people, while Algeria was projected
to reach approximately 51.20 million people.

In 2022, China was the most populous country with
1.426 billion people, followed by India with
1.417 billion. The United States ranked third
with approximately 338.29 million people.
Population trend visualizations showed changes
in population among the ten selected countries
between 1990 and 2022.

Correlation analysis revealed a moderate positive
relationship between land area and world population
percentage (r=0.4533). Other relationships were
weak. The correlation between 2022 population size
and 2019 CO2 emissions per capita was
approximately r=0.00817, indicating almost no
linear association.

The continental comparison showed that Asia had
the highest sum of country-level per-capita
emission values (242), followed by Europe (229).
North America ranked third (79.6), while South
America recorded the lowest value (29.4). These
values represent sums of per-capita emissions,
not actual continental total emissions.


\section{Discussion}

The findings indicate that population growth varies
 across countries, with China and India
having the largest populations in 2022. The population
projections of 2030  demonstrate how differences in
growth rates can influence future population sizes.

The correlation coefficient of 0.00817 suggests
that population size has almost no linear relationship
with CO2 emissions per capita in the merged
dataset (merged\_data). This is indicated by
energy consumption, industrial activities and
economic development among countries.

The comparison among continents showed that Asia had
the highest sum of country-level per-capita
emissions, followed by Europe. However, these sums
should not be interpreted as total continental
emissions because the analysis used per-capita
values rather than total emissions.



\section{Lessons Learned and Recommendations}

This project improved my level of  understanding more about data
imprtation, data management, exploratory data analysis, data
visualization and correlation analysis using R.
I became more familiar with using R programming 
commands,learnt how mathematical models
can be applied to project future population sizes.

For population trends, learnt the relation between
 World population and CO2 emissions,the impact CO2 emissions
have on the environment . Additional variables,
such as GDP, industrialization and energy
consumption.Furthermore, population projections should be
interpreted cautiously because growth rates may
change over time. And here are some recommended links 
that were used to help in analysing our worls Population and CO2 trends 
//https://www.kaggle.com/code/devraai/world-gdp-
population-co2-analysis?
{https://https://www.kaggle.com/datasets/iamsouravbanerjee
/world-population-dataset?



\section{Conclusion}

This study examined world population trends and
their relationship with CO2 emissions using
statistical analysis and visualization in R.
The results identified China and India as the
most populous countries in 2022 and highlighted
differences in population growth across countries.
The correlation analysis revealed almost no
linear relationship between population size and
per-capita CO2 emissions. Overall, the study
demonstrates the importance of combining
population, economic and environmental indicators
to better understand global emission patterns.





\end{document}

